\section{1D Non-Vertical Scenario}

Here, we will assume a 1D scenario (that is, we assume that the velocity varies only with height, being constant in the other two dimensions), with a single source at the bottom, and an array of $N$ receivers placed vertically along the water column; however, we don't necessarily require a vertical alignment between the source and the receivers, where these can be horizontally displaced.

We also assume an ideal noiseless environment, where each measurement reflects exactly the wave's travel time, and the positions for all devices are precisely known.

\subsection{Discrete method}
\label{subsec:sec3:discrete_method}

\textit{This approach can be seen as the solution to the Eikonal equation, as detailed in \cite{ali_opensource_2019}, but doing so with only one dimension. Their process works iteratively over all media, while here we propose solving each layer \quote{pixel} at a time.}

This is one of the possible ways for solving the SSP estimation problem, with TTT. The body is discretized along the varying dimension(s), with receivers positioned along it. The speed of sound is assumed to be constant within each cell, and thus the travel-time will be the summation of the distance traveled in the cell, divided by the speed in the cell. This, done for each sensor/measurement, gives a system of equations which is used to estimate the speed in each cell.

Given enough measurements (in this case, enough receivers), this technique can have cells of an arbitrary size, therefore achieving an arbitrary accuracy. Among its advantages are its simplicity, robustness, and computational efficiency. However, it depends on the accurate knowledge on the position of all devices (sources and receivers), as well as precise timing between all of them.

Along with the sources and receivers positions, we also assume an ideal noiseless environment, where each measurement reflects exactly the wave's travel time. \Cref{fig:sec2:wavepaths_heterogeneous-variable_environment:fig1} shows the path of a wave produced by the source (in green) to each of the receivers (in blue).

\begin{figure}[H]
	\centering
	\begin{tikzpicture}[scale=1.5]
		\coordinate (tl) at (-3, 2);
		\coordinate (tr) at (3, 2);
		\coordinate (bl) at (-3, -2);
		\coordinate (br) at (3, -2);
		\coordinate (cc) at (0,0);
		\coordinate (tc) at (0,2);
		\coordinate (bc) at (0,-2);
		\coordinate (cl) at (-3, 0);
		\coordinate (cr) at (3,0);
		\coordinate (source) at ($(bl)!0.95!(br)$);
		\coordinate (receiver) at ($(bl)!0.9!(tl)$);
		\draw[gray,fill=none] (tl) rectangle (br);
		\foreach \x/\val in {0.8/$r_1$, 0.6/$r_2$, 0.4/$r_3$, 0.2/$\vdots\,$, 0/$r_N$}
		{
			\pgfmathparse{\x+0.2} \let\y=\pgfmathresult 
			\coordinate (pl) at ($(tl)!\y!(bl)$); %Point-left
			\coordinate (ppl) at ($(tl)!\x!(bl)$);
			\coordinate (ps) at ($(pl)!0.5!(ppl)$);
			\pgfmathparse{(0.8-\x)*-12/0.8 -6} \let\oang=\pgfmathresult
			\pgfmathparse{(0.8-\x)*-54/0.8 + 174} \let\iang=\pgfmathresult
			\draw [->-,thick] (source) to[out=\iang,in=\oang] (ps);
			\node[square,fill=blue!70!black] () at (ps) {};
			\node[left=6pt of ps] () {\val};
		}
		\node[triangle,fill=green!70!black] () at (source) {};
		\node[below=6pt of source] () {$s$};
		
	\end{tikzpicture}
	\caption{Wave paths on heterogeneous and continuously variable environment. The green triangle represents a source, and the blue squares each of the $N$ receivers.}
	\label{fig:sec2:wavepaths_heterogeneous-variable_environment:fig1}
\end{figure}

We now consider that each receiver represents a region in the medium (in this case, a layer), with the receiver being positioned at the center of its layer; and model the system such that each stratum has a constant speed. This will result in \cref{fig:sec2:wavepaths_heterogeneous-variable_environment:fig2}, where the wave now discontinuously changes direction and speed when changing media; however, its wavepath is continuous. The different layers are represented by the different shades of gray.

\begin{figure}[H]
	\centering
	\begin{tikzpicture}[scale=1.5]
		\coordinate (tl) at (-3, 2);
		\coordinate (tr) at (3, 2);
		\coordinate (bl) at (-3, -2);
		\coordinate (br) at (3, -2);
		\coordinate (cc) at (0,0);
		\coordinate (tc) at (0,2);
		\coordinate (bc) at (0,-2);
		\coordinate (cl) at (-3, 0);
		\coordinate (cr) at (3,0);
		\coordinate (source) at ($(bl)!0.95!(br)$);
		\coordinate (receiver) at ($(bl)!0.9!(tl)$);
		\draw[gray,fill=none] (tl) rectangle (br);
		\foreach \x in {0.8,0.6,0.4,0.2,0}
		{
			\coordinate (pr) at ($(tr)!\x!(br)$); % Point-right
			\pgfmathparse{\x+0.2} \let\y=\pgfmathresult 
			\coordinate (pl) at ($(tl)!\y!(bl)$); %Point-left
			\pgfmathparse{\x * 50} \let\y=\pgfmathresult % Color calculation
			\draw[gray,very thin,fill=gray!\y!white] (pl) rectangle (pr);
		}
		\draw [->-,thick] (source) to (-3.0,-1.6);
		%
		\draw [->-,thick] (source) to (-0.4,-1.2) to (-3.0,-0.8);
		%
		\draw [->-,thick] (source) to (+1.25,-1.2) to (-1.14,-0.4) to (-3.0,0.0);
		%		
		\draw [->-,thick] (source) to (+1.83,-1.2) to (0.6,-0.4) to (-1.53,0.4) to (-3.0,0.8);
		%		
		\draw [->-,thick] (source) to (+2.15,-1.2) to (1.4,-0.4) to (0.2,0.4) to (-1.76,1.2) to (-3.0,1.6);
		%
		\foreach \x/\val in {0.8/$r_1$, 0.6/$r_2$, 0.4/$r_3$, 0.2/$\vdots\,$, 0/$r_N$}
		{
			\pgfmathparse{\x+0.2} \let\y=\pgfmathresult 
			\coordinate (pl) at ($(tl)!\y!(bl)$); %Point-left
			\coordinate (ppl) at ($(tl)!\x!(bl)$);
			\coordinate (ps) at ($(pl)!0.5!(ppl)$);
			\node[square,fill=blue!70!black] () at (ps) {};
			\node[left=6pt of ps] () {\val};
		}
		\node[triangle,fill=green!70!black] () at (source) {};
		\node[below=6pt of source] () {$s$};
		
	\end{tikzpicture}
	\caption{Wave paths on heterogeneous and discretized environment.}
	\label{fig:sec2:wavepaths_heterogeneous-variable_environment:fig2}
\end{figure}

\subsubsection{First layer}
The speed on the bottom-most layer can be trivially calculated. With the precise position of all sensors as well as for the source, 
\begin{equation}
	c_1 = \frac{d_1}{t_1},
\end{equation}
where $c_1$ is the speed on that layer; $d_1$ is the distance between the source and the first receiver; and $t_1$ is the travel time between those two points.

\subsubsection{Second layer}
Considering now the next layer and the wave that hits its respective sensor, we have the following scheme (with the $\y$-axis scaled, for clarity):

\begin{figure}[H]
	\centering
	\begin{tikzpicture}[x=1.5cm,y=3.5cm]
		\coordinate (tl) at (-3, 2);
		\coordinate (tr) at (3, 2);
		\coordinate (bl) at (-3, -2);
		\coordinate (br) at (3, -2);
		\coordinate (cc) at (0,0);
		\coordinate (tc) at (0,6);
		\coordinate (bc) at (0,-2);
		\coordinate (cl) at (-3, 0);
		\coordinate (cr) at (3,0);
		%
		\coordinate (source) at ($(bl)!0.95!(br)$);
		\coordinate (r1) at ($(bl)!0.1!(tl)$);
		\coordinate (r2) at ($(bl)!0.3!(tl)$);
		%
		\coordinate (p0) at ($(tr)!0.6!(br)$);
		\coordinate (p1) at ($(tr)!0.8!(br)$);
		\coordinate (ti) at (-0.4,-0.4); % top interface
		\coordinate (bi) at (-0.4,-2); % bottom interface
		\coordinate (ci) at (-0.4,-1.2); % center interface
		%
		\draw[gray,fill=none] ($(tl)!0.6!(bl)$) rectangle (br);
		%% Draw colored rectangles
		\foreach \x in {0.8, 0.6}
		{
			\coordinate (pr) at ($(tr)!\x!(br)$); % Point-right
			\pgfmathparse{\x+0.2} \let\y=\pgfmathresult 
			\coordinate (pl) at ($(tl)!\y!(bl)$); %Point-left
			\pgfmathparse{\x * 50} \let\y=\pgfmathresult % Color calculation
			\draw[gray,very thin,fill=gray!\y!white] (pl) rectangle (pr);
		}
		%% Draw wave paths
		%		\draw [->-,thick] (source) to (r1);
		%
		\draw [->-,thick] (source) to (ci) to (r2);
		%% Draw receivers and sources
		\node[square,fill=blue!70!black] () at (r2) {};
		\node[left=6pt of r2] () {$r_2$};
		%
		\node[triangle,fill=green!70!black] () at (source) {};
		\node[below=6pt of source] () {$s$};
		%% Draw distances
		\draw[<->,thin,gray]
		([xshift=0.2cm]p1) -- ([xshift=0.2cm]br) node[pos=0.5,anchor=west,black] {$h_1$};
		%
		\draw[<->,thin,gray]
		([xshift=0.2cm]p0) -- ([xshift=0.2cm]p1) node[pos=0.5,anchor=west,black] {$h_2$};
		%
		\draw[<->,thin,gray]
		([xshift=0.15cm,yshift=0.15cm]ci) -- ([xshift=0.15cm,,yshift=0.15cm]source)
		node[pos=0.5,anchor=south west,black] {$d_{2,1}$};		
		%
		\draw[<->,thin,gray]
		([xshift=-0.15cm,yshift=-0.15cm]ci) -- ([xshift=-0.15cm,,yshift=-0.15cm]r2)
		node[pos=0.2,anchor=north east,black] {$d_{2,2}$};		
		%
		\draw[<->,thin,gray] ([xshift=-0.25cm,yshift=-0.6cm]source) -- ([xshift=-0.25cm,yshift=-0.6cm]r2) node[pos=0.4, anchor=north east,black] {$d_{2}$};
		%% Draw angles
		\draw[dashed, semithick]  (bi) -- (ti);
		\draw pic[draw,angle radius=0.5cm, "$\theta_1$", angle eccentricity=1.8] {angle=bi--ci--source};
		%
		\draw pic[draw,angle radius=0.5cm, "$\theta_2$", angle eccentricity=1.8] {angle=ti--ci--r2};
		
	\end{tikzpicture}
	\caption{Wave path in a heterogeneous and discrete environment, for the second receiver.}
	\label{fig:sec2:wavepaths_heterogeneous-variable_environment:fig3}
\end{figure}

From Snell's law, we have that
\begin{equation}
	\frac{c_1}{c_2} = \frac{\sin \theta_1}{\sin \theta_2},
\end{equation}
where $\theta_n$ is the wave's direction-of-travel's angle -- which will be called wavepath's angle -- w.r.t. the vector normal to the separation between the layers. This leads to
\begin{equation}
	c_2 = \frac{c_1 \sin \theta_2}{\sin \theta_1}
\end{equation}

The travel time for the wave that hits the second sensor will be
\begin{equation}
	t_2 = \frac{d_{2,1}}{c_1} + \frac{d_{2,2}}{c_2},
\end{equation}
in which $d_{2,1}$ is the distance traveled by the wave arriving at the second sensor, while in the first (bottom-most) layer; and $d_{2,2}$ is the distance traveled while in the second layer.

Geometrically, we also have that
\begin{equation}
	\begin{split}
		h_1
		& = \cos\theta_1 d_{2,1} \\
		& = \Sqrt{1-\sin^2\theta_1} d_{2,1}
	\end{split},
\end{equation}
\begin{equation}
	\frac{h_2}{2} = \Sqrt{1-\frac{c_2^2}{c_1^2}\sin^2\theta_1}d_{2,2}\,,
\end{equation}
where $h_1$ and $h_2$ are the heights of each receiver's layers. Defining $\sigma_1 = \sin\theta_1$, we then have
\begin{equation}
	t_2 = \frac{h_1}{c_1 \Sqrt{1-\sigma_1^2}} + \frac{h_2}{2c_2\Sqrt{1-\frac{c_2^2}{c_1^2} \sigma_1^2}}.
	\label{eq:sec3:ttt:travel-time_t2}
\end{equation}

Letting $\tau_1 = \tan \theta_1$, geometrically we have that
\begin{equation}
	d_2^2 = \pts{h_1\tau_1 + \frac{h_2}{2}\tau_2}^2 + \pts{h_1 + \frac{h_2}{2}}^2,
\end{equation}
which, via Snell's law, and trigonometrical properties, become
\begin{equation}
	d_2^2 = \pts{h_1 \frac{\sigma_1}{\Sqrt{1-\sigma_1^2}} + \frac{h_2}{2} \frac{\frac{c_2}{c_1}\sigma_1}{\Sqrt{1 - \frac{c_2^2}{c_1^2} \sigma_1^2}} }^2 + \pts{h_1 + \frac{h_2}{2}}^2,
	\label{eq:sec3:ttt:distance_d2}
\end{equation}
with $d_2$ being the distance between the source and the second (second bottom-most) receiver.

Given that $h_1$, $h_2$, $t_2$, $c_1$ and $d_2$ are known values, this results in a system with two non-linear equations \cref{eq:sec3:ttt:distance_d2,eq:sec3:ttt:travel-time_t2}, and two unknowns $\sigma_1$ and $c_2$. Solving this system leads to the knowledge of $c_2$.

%To find this solution, the Newton-Raphson method can be used. Let $f(\bvs)$ and $g(\bvs)$ be defined as
%\begin{gather*}
%	f(\bvs) = \pts{h_1 \frac{\sigma}{\Sqrt{1-\sigma^2}} + \frac{h_2}{2} \frac{\frac{c}{c_1}\sigma}{\Sqrt{1 - \frac{c^2}{c_1^2} \sigma^2}} }^2 + \pts{h_1 + \frac{h_2}{2}}^2 - d_2^2 \\
%	g(\bvs) = \frac{h_1}{c_1 \Sqrt{1-\sigma^2}} + \frac{h_2}{2c\Sqrt{1-\frac{c^2}{c_1^2} \sigma^2}} - t_2
%\end{gather*}
%with $\bvs = \tr{[c,~\sigma]}$; $c$ representing the velocity in the second layer, and $\sigma$ the emission angle of the wave ray that hits the second sensor. We define $\opt{\bvs}$ as the zero of this system of equations, in which $f(\opt{\bvs}) = g(\opt{\bvs}) = 0$.
%
%Let $F(\bvs)$ be a vector of the two previous functions, and $J(\bvs)$ its Jacobian, both given by
%\begin{gather*}
%	F(\bvs) = \begin{bmatrix}
	%		f \\
	%		g
	%	\end{bmatrix} (\bvs)\\
%	J(\bvs) = \begin{bmatrix}
	%		\pdv{f}{c} & \pdv{f}{\sigma} \\
	%		\pdv{g}{c} & \pdv{g}{\sigma}
	%	\end{bmatrix}(\bvs)
%\end{gather*}
%
%With this, and with an initial guess $\bvs{0}$, we define the iterative process
%\begin{equation}
%	\bvs{k+1} = \bvs{k} - \inv{J(\bvs)} \ast F(\bvs)
%\end{equation}
%through which we can achieve arbitrary precision; that is, we can iteratively approach $\opt{\bvs}$, leading to an arbitrarily small $F(\bvs{k+1})$.

\subsubsection{Generic layer}

We now consider the $m$-th layer, and assume that all $m-1$ previous layers were solved, knowing the velocities $c_1$ through $c_{m-1}$.

Considering the $m$-th layer, we have that
\begin{equation}
	t_m = \frac{h_m}{2c_m\Sqrt{1-\frac{c_m^2}{c_1^2}\sigma_{m,1}^2}} + \sum_{i=1}^{m-1} \frac{h_i}{c_i\Sqrt{1-\frac{c_i^2}{c_1^2}\sigma_{m,1}^2}}
	\label{eq:sec3:ttt:travel-time_mth-layer}
\end{equation}
and
\begin{equation}
	d_m^2 = \pts{\frac{h_m}{2} \frac{\frac{c_m}{c_1}\sigma_{m,1}}{\Sqrt{1-\frac{c_m^2}{c_1^2}\sigma_{m,1}^2 }} + \sum_{i=1}^{m-1} h_i \frac{\frac{c_i}{c_1}\sigma_{m,1}}{\Sqrt{1-\frac{c_i^2}{c_1^2}\sigma_{m,1}^2 }} }^2 + \pts{\frac{h_m}{2} + \sum_{i=1}^{m-1} h_m}^2.
	\label{eq:sec3:ttt:distance_mth-layer}
\end{equation}

This leads to a system with two non-linear equations \cref{eq:sec3:ttt:travel-time_mth-layer,eq:sec3:ttt:distance_mth-layer}, and two unknowns $c_m$ and $\sigma_{m,1}$; where $\sigma_{m,1} = \sin\theta_{m,1}$ is the $m$-th sensor's wavepath direction w.r.t. the normal, in the first layer. This can be solved in an identical way as for the second layer case, using the appropriate functions and variables.

Note that, although the previous $m-1$ layers' velocities are necessary to solve the $m$-th one, this process can be sped by doing the iterative search simultaneously for all velocities. That is, instead of achieving a sufficiently precise velocity for $c_2$, then $c_3$, repeating until $c_m$ is found, the iteration for all velocities can be done in the same loop, and the search for one only stop if all previous velocities have achieved sufficient precision.

%\subsubsection{Noisy measurements}
%
%We now assume that there is imprecision in our measurements, for both devices' position, and travel time. This could be caused by the devices' movements (either sources or receivers), synchronization mismatches between the devices, and unknown interactions (reflections, refractions, etc) between the waves and the environment.

%\section{Interpolation method}
%
%In this technique, we assume that the body of water follows a polynomial curve along the column of water. Similarly to the discrete method, here we assume that the speed is constant horizontally, and varies vertically. However, now we take the speed to continuously change with the height, instead of discretely.
%
%Assuming 

\subsection{Continuous fitting}

Consider any point $\bvp$ along the curve $\bvr[v]$. According to Snell's law, it will follow the equation
\begin{equation}
	\frac{\sigma(\bvp)}{c(\bvp)} = \frac{\sigma_s}{c_s},
\end{equation}
where $\sigma_s$ and $c_s$ are the (sine of) the angle of the wavepath at the source, and $c_s$ is the wave speed at the source.

Assuming that $\bvp$ is on the curve $\bvr[v]$ and is a function of time $t$, its derivative (that is, the variation of its position along $\bvr[v]$ over time) will be given by
\begin{equation}
	\dv{\bvp}{t} = [c_{\x},~c_{\y}],
\end{equation}
with $c_{\x}$ and $c_{\y}$ being the components of the wave's speed along each direction. Geometrically, this can be rewritten as
\begin{equation}
	\dv{\bvp}{t} = [c(\bvp) \sigma(\bvp),~c(\bvp) \gamma(\bvp)],
\end{equation}
where $\gamma(\bvp) = \Sqrt{1-\sigma^2(\bvp)}$, or the cosine of the angle of the wave's direction at $bvp[v]$. Through Snell's law, we then have that
\begin{equation}
	\dv{\bvp}{t} = \bts{\frac{\sigma_s}{c_s}c^2(\bvp),~c(\bvp)\Sqrt{1-\frac{\sigma_s^2}{c_s^2}c^2(\bvp)}}.
\end{equation}

This forms a system of non-linear differential equations. Given the velocity profile $c(\bvx)$ (where $\bvx$ represents a position generically in the body of water), and initial point -- the source -- $\bvp{s}$ with velocity $c_s$ and angle $\theta_s$, it is possible to determine the path traced by this wave. This is the direct problem, to determine the path given the SSP.

Meanwhile, in the inversion problem at hand there is a source, a set of arrival points $r_1$ through $r_N$ with travel-time measurements $t_1$ through $t_N$, and the objective is reconstructing the SSP $c(\bvx)$.

\subsubsection{Polynomial fitting}

The simplest solution for this problem is through a polynomial fitting. Having $N$ travel-time measurements, an $(N-1)$-degree polynomial can be fitted along these points.

We assume that $c(\bvp) \equiv c(y)$ is constant in the $\x$-axis, varying only on the $\y$-axis. It can then be written as
\begin{equation}
	c(y) = \sum_{n=0}^{N-1} a_n y^n.
\end{equation}

Writing $\bvp(t) = [p_x(t),~p_y(t)]$, then
\begin{gather}
	\dv{p_x}{t} = \frac{\sigma_s}{c_s} c^2\pts{p_y} \\
	\dv{p_y}{t} = c\pts{p_y} \Sqrt{1-\frac{\sigma_s^2}{c_s^2}c^2\pts{p_y}}.
\end{gather}

%Assuming that the right hand-side of the equations are constants w.r.t. time -- that is, assuming that the SSP $v(y)$ is constant over time, then
%\begin{gather}
%	p_x(t) = \pts{\frac{\sigma_0}{c_0} v^2(y)} t + p_{x;0} \\
%	p_y(t) = \frac{\sigma_0}{c_0} v(y) \Sqrt{1-\frac{\sigma_0^2}{c_0^2}v^2(y)} t + p_{y;0}
%\end{gather}

Note that the second equation doesn't depend on $x$, being an ordinary differential equation, with the first one depending on its solution.

We let $\lambda_m = \dfrac{\sigma_{m,s}}{c_s}$ be the wavepath's angle -- at source -- of the wavepath that hits the $m$-th sensor, and $p_{y,m}(t)$ the vertical component of the wavepath that reaches the $m$-th sensor; and via separation of variables,
\begin{equation}
	\frac{\dd p_{y,m}}{c\pts{p_{y,m}}\Sqrt{1-\lambda_1^2 c^2\pts{p_{y,m}}}} = \dd t.
\end{equation}

Integrating both sides, we achieve
\begin{equation}
	\int \frac{1}{c\pts{p_{y,m}}\Sqrt{1-\lambda_1^2 c^2\pts{p_{y,m}}}} \dd p_{y,m} = t + t_0,
\end{equation}
which, even with the knowledge/assumption that $c\pts{p_{y,m}}$ is a polynomial, still doesn't necessarily have a closed-form solution.

Together with the previous equation, which states
\begin{equation}
	\dv{p_{x,m}}{t} = \lambda_1 c^2\pts{p_{y,m}},
\end{equation}
which can also be integrated, leading to
\begin{equation}
	p_{x,m} + p_0 = \int \lambda_1 c^2\pts{p_{y,m}(t)} \dd t,
\end{equation}
where $p_{y,m}(t)$ is the vertical position as a function of time. This brings another problem, that $p_{y,m}(t)$ isn't directly obtainable, given that the previous results lead to $t(p_{y,m})$; that is, the time of arrival of the wave-path, at a given height. However, from the geometry of the problem we also have that $p_{x,m}(t)$ and $p_{y,m}(t)$ are monotonously functions. Therefore, $t(p_{y,m})$ can be inverted numerically, at least within the region of space considered in the problem.

With this, we have a system with $2N$ equations ($N$ relating to the position in the $\y$-axis, and $N$ to the $\x$-axis), and $2N$ variables (the $N$ parameters that control the polynomial $c(y)$, and the $N$ wavepath angles $\theta_{m,s}$, from which $\sigma_{m,s} = \sin \theta_{m,s}$).

Therefore, a solution to the problem can be achieved numerically and iteratively, searching for the optimal parameters. Starting with initial parameters for all variables, these can be updated iteratively, until sufficient precision is achieved.

\subsubsection{Polynomial inverse fitting}

Here, instead of assuming that $c(y)$ has a polynomial behavior, we assume that $s(y) = \frac{1}{c(y)}$ -- the slowness -- is polynomial. That is,
\begin{equation}
	s(y) = \frac{1}{c(y)} = \sum_{n=0}^{N-1} b_n y^n.
	\label{eq:sec2:polynomial_slowness}
\end{equation}

With this, the previous modeling becomes
\begin{equation}
	\begin{split}
		\dd t
		& = \frac{\dd p_{y,m}}{\frac{1}{s(p_{y,m})} \Sqrt{1-\lambda_1^2\frac{1}{s^2(p_{y,m})}}}\\
		& = \frac{s^2(p_{y,m})}{\Sqrt{s^2(p_{y,m}) - \lambda_1^2}} \, \dd p_{y,m}.
	\end{split},
\end{equation}

This modeling's solving also doesn't have a closed form, and would require a numerical-iterative solution.
